\scanpage{91}
\uline{NB} Si on a un hom
\[
u:\underline{B}\longrightarrow\underline{\mathrm{Pic}}_{X/S}
\]
avec $B$ un schéma abélien, et si $\underline{\Omega}^1_{X/S}$ \unsure{coh.|est} plat/$S$ en dim 0, alors $u$ se factorise en $B\longrightarrow\underline{\mathrm{Pic}}^\omega_{X/S}$. En particulier, si $\underline{\mathrm{Pic}}^{00}_{X/S}$ existe, alors on a
\[
\underline{\mathrm{Pic}}^{00}_{X/S}\subset\underline{\mathrm{Pic}}^\omega_{X/S}.
\]
Notons que la \uline{donnée} de $u$ équivaut : à \unsure{la} donnée d'un torseur $A$ sous \deleted{\unsure{A}} le schéma dual $B^*$ de $B$, et d'un morphisme
\[
X\longrightarrow A
\]
\margin{sous réserve d'existence de $B^*$ et de \unsure{possibilité} \unsure{en torseurs} (p.ex. si $B/S$ est projectif)}
\[
[\text{d'où}\qquad B\simeq\underline{\mathrm{Pic}}^0_{A/S}\subset\underline{\mathrm{Pic}}_{A/S}\longrightarrow\underline{\mathrm{Pic}}_{X/S}\quad\ldots\ ]
\]
On trouve alors, par fonctorialité, un diagramme
\margin{fonctorialité doit être explicitée avec un peu de soin !}
\diagram{\begin{tikzcd}[column sep=large,row sep=large]
0 \arrow[r,dash] & \underline{\omega}^{\unsure{*|+}}_{X/S} \arrow[r] & M_{X/S} \arrow[r] & \underline{\mathrm{Pic}}^\omega_{X/S} \arrow[r] & 0 \\
0 \arrow[r] & \underline{\omega}_{A/S} \arrow[r]\arrow[u] & M_{A/S} \arrow[r]\arrow[u] & \underline{\mathrm{Pic}}^\omega_{A/S} \arrow[r,dash]\arrow[u] & 0 \\
0 \arrow[r,dash] & \underline{\omega}_{A/S} \arrow[r]\arrow[u,equal] & \deleted{\unsure{E_{B/S}}}\,M^0_{A/S} \arrow[r]\arrow[u] & B \arrow[r]\arrow[u,hook] & 0 \\
& \mathrm{V}(\mathrm{R}^1\unsure{f_{B*}}(\mathcal{O}_B)) \arrow[u,"\simeq"{sloped},draw=none] & & \underline{\mathrm{Pic}}^0_{A/S} \arrow[u,equal] &
\end{tikzcd}}
Ici $\deleted{\unsure{E_{B/S}}}\inserted{M^0_{A/S}}$ est défini comme image inverse de $M_{A/S}$ \inserted{par $B\longrightarrow\underline{\mathrm{Pic}}^\omega_{A/S}$} [ou aussi de $M_{X/S}$ par $B\longrightarrow\underline{\mathrm{Pic}}^\omega_{X/S}$ \inserted{si $\underline{\omega}_{A/S}\simeq\underline{\omega}_{X/S}$}]. Je vou-
\margin{90}
